\section{Experiments and Analysis}
\label{sec:experiments}

\subsection{Comparison with Frontier Models}
\label{sec:exp:baselines}

To contextualize our fine-tuned 8B model's performance, we evaluate three frontier models on two benchmarks: (1) a 500-question sample from our training dataset (v4), and (2) a fully independent 54-question holdout set composed of textbook problems, professor-authored exam questions, and French reliability exercises (\emph{TD fiabilit\'{e}}), none of which appear in the training data. For all models, we use the same answer extraction and comparison pipeline described in Section~\ref{sec:sft:eval}. When the extraction pipeline fails to parse a model's response (\textsc{unable\_to\_extract}), we exclude that question from the denominator rather than counting it as incorrect, to avoid penalizing models for formatting differences.

\begin{table}[t]
\centering
\small
\caption{Frontier model accuracy on our reliability engineering benchmarks. $N_{\text{eval}}$ is the number of questions with successfully extracted answers.}
\label{tab:baselines_v4}
\begin{tabular}{lccc}
\toprule
\multirow{2}{*}{Model} & \multicolumn{3}{c}{v4 500-sample} \\
\cmidrule(lr){2-4}
 & Correct & $N_{\text{eval}}$ & Accuracy \\
\midrule
Claude Sonnet 4.6      & 462 & 500 & \textbf{92.4\%} \\
o3-mini                & 427 & 494 & 86.4\% \\
Gemini 3.1 Pro Preview & 308 & 403 & 76.4\% \\
\midrule
Qwen3-8B (base)        & \multicolumn{3}{c}{63.6\% (10-fold CV)} \\
Qwen3-8B (SFT only)    & \multicolumn{3}{c}{82.2\% (10-fold CV)} \\
\textbf{Qwen3-8B (SFT+GRPO)} & \multicolumn{3}{c}{\textbf{57.9\%} on 266 hard q.} \\
\bottomrule
\end{tabular}
\end{table}

\begin{table}[t]
\centering
\small
\caption{Independent holdout evaluation (54 questions, no overlap with training data).}
\label{tab:baselines_holdout}
\begin{tabular}{lccc}
\toprule
\multirow{2}{*}{Model} & \multicolumn{3}{c}{Holdout (54 questions)} \\
\cmidrule(lr){2-4}
 & Correct & $N_{\text{eval}}$ & Accuracy \\
\midrule
Claude Sonnet 4.6      & 43 & 53 & \textbf{81.1\%} \\
o3-mini                & 37 & 48 & 77.1\% \\
Gemini 3.1 Pro Preview & 19 & 23 & 82.6\% \\
\midrule
Qwen3-8B (base)        & 29 & 54 & 53.7\% \\
Qwen3-8B (SFT fold\_4) & 18 & 54 & 33.3\% \\
Qwen3-8B (SFT+GRPO)   & 29 & 54 & 53.7\% \\
\bottomrule
\end{tabular}
\end{table}

On the v4 sample (Table~\ref{tab:baselines_v4}), Claude Sonnet 4.6 leads at 92.4\%, followed by o3-mini (86.4\%) and Gemini 3.1 Pro (76.4\%). Our fine-tuned Qwen3-8B achieves 82.2\% under cross-validation on the same dataset, placing it between o3-mini and Gemini 3.1 Pro despite being an order of magnitude smaller. Note that the fine-tuned model's accuracy is measured via 10-fold CV (training on 90\%, testing on 10\%) and is therefore not directly comparable to the frontier models which see all questions at test time; the comparison is nevertheless informative as an approximate ranking.

On the independent holdout (Table~\ref{tab:baselines_holdout}), Claude Sonnet 4.6 achieves 81.1\% and o3-mini 77.1\%. Gemini 3.1 Pro achieves 82.6\% on the 23 questions where answer extraction succeeded, though the high extraction failure rate (31/54 = 57.4\%) limits the reliability of this estimate. Our base Qwen3-8B achieves 53.7\%, but the SFT model regresses sharply to 33.3\% ($-20.4$pp), revealing severe catastrophic forgetting on truly unseen questions. GRPO fully recovers to 53.7\%, matching the base model and demonstrating that RL training restores generalization lost during SFT while adding in-distribution gains (Section~\ref{sec:grpo:results}).

\subsection{General Reasoning Evaluation}
\label{sec:exp:reasoning}

A common concern with domain-specific fine-tuning is catastrophic forgetting of general capabilities. To assess whether our SFT and GRPO pipeline preserves or improves general reasoning, we evaluate the base and fine-tuned models on standard mathematical reasoning benchmarks.

Our holdout evaluation (Table~\ref{tab:baselines_holdout}) provides direct evidence of catastrophic forgetting: SFT fold\_4 drops from 53.7\% to 33.3\% on 54 unseen reliability engineering questions. This is consistent with the observation that SFT was trained exclusively on 600 ``easy'' questions (where the base model already scores 4/4), which teaches surface patterns without improving hard-problem reasoning. GRPO partially recovers this regression (50.0\%), suggesting that RL training restores some of the base model's general reasoning capability lost during SFT.

\subsection{Analysis}
\label{sec:exp:analysis}

\subsubsection{Dataset Size Scaling}

Figure~\ref{fig:scaling} plots SFT improvement as a function of training set size. The relationship is clearly non-linear: improvements are marginal below 300 examples ($+2.0$--$2.3\%$), begin to emerge at 501 examples ($+6.2\%$, $p=0.063$), and become large and statistically significant at 866 examples ($+18.6\%$, $p=0.002$). This suggests a threshold effect around 500 examples, below which LoRA fine-tuning on this domain struggles to overcome the regularization imposed by the frozen base weights.

\TODO{Add scaling curve figure.}

\begin{table}[t]
\centering
\small
\caption{SFT improvement and question flips across dataset sizes. W$\to$R: wrong$\to$right, R$\to$W: right$\to$wrong.}
\label{tab:scaling}
\begin{tabular}{rccccc}
\toprule
$N$ & $\Delta$ & $p$ & W$\to$R & R$\to$W & Ratio \\
\midrule
215 & $+2.3$ & 0.50  & 23 & 18 & 1.3 \\
280 & $+2.0$ & 0.69  & 19 & 15 & 1.3 \\
501 & $+6.2$ & 0.063 & 74 & 43 & 1.7 \\
\textbf{866} & $\boldsymbol{+18.6}$ & \textbf{0.002} & \textbf{206} & \textbf{45} & \textbf{4.6} \\
\bottomrule
\end{tabular}
\end{table}

The question flip ratio (Table~\ref{tab:scaling}) is particularly informative: at 215 samples, the model trades nearly as many correct answers as it gains (1.3:1), while at 866 samples, gains overwhelm regressions (4.6:1). This indicates that larger augmented datasets not only improve average performance but also make the improvement more \emph{robust}, with fewer catastrophic regressions on individual questions.

\subsubsection{Epoch Convergence}

Table~\ref{tab:epochs} shows SFT improvement as a function of training epochs on the 866-question dataset (10-fold CV). Performance increases rapidly from 3 to 4 epochs, then plateaus. Early stopping with patience 2 independently selects epoch 4 in both Exp.~22 (max 6) and Exp.~24 (max 8), confirming this as a stable convergence point.

\begin{table}[t]
\centering
\small
\caption{Effect of training epochs on 866-question dataset (10-fold CV).}
\label{tab:epochs}
\begin{tabular}{cccc}
\toprule
Epochs & FT Accuracy & $\Delta$ & $p$ \\
\midrule
3            & 75.1\% & $+11.4$ & 0.002 \\
4            & 79.3\% & $+15.7$ & 0.002 \\
6 (ES$\to$4) & 81.8\% & $+18.1$ & 0.002 \\
8 (ES$\to$4) & 82.2\% & $+18.6$ & 0.002 \\
\bottomrule
\end{tabular}
\end{table}

\subsubsection{Error Analysis}

A detailed error analysis by sub-topic is left for future work. We note that the GRPO model's truncation rate increases with training (4 for base, 10 for ckpt-80), suggesting that GRPO encourages longer reasoning chains that occasionally exceed the generation limit. The holdout accuracy matches the base model exactly (53.7\%), suggesting that GRPO improves in-distribution performance without degrading general reasoning---unlike SFT alone, which causes severe holdout regression.
